% =====================================================================
%  2bat-ud6-13.tex  —  SESSIÓ 13: Síntesi de l'etapa: la mirada crítica amb dades
%  (sense preàmbul: s'inclou des de main.tex)
% =====================================================================
\horatitol{Sessió 13 — Síntesi de l'etapa: la mirada crítica amb dades}%
{Traçar el recorregut estadístic complet del batxillerat —descriure, modelar, inferir i criticar— i reconèixer l'habilitat transferible que en resulta.}

\subsection{Idea clau}

Penúltima sessió de contingut de tota l'etapa. Fem una \textbf{síntesi} del recorregut
estadístic complet del batxillerat professionalitzador, que és la columna vertebral de la
matèria per a aquest perfil. Tot aquest itinerari condueix a una única habilitat
transferible: \textbf{saber llegir, interpretar i qüestionar dades sobre persones i
col·lectius}.

\begin{idea}
\textbf{El recorregut estadístic de l'etapa, en quatre passos}
\begin{enumerate}[leftmargin=1.6em, itemsep=2pt]
  \item \textbf{Descriure} (1r): estadística descriptiva —paràmetres, gràfics— per resumir
        dades.
  \item \textbf{Mesurar la incertesa} (1r): probabilitat —la possibilitat d'un cas.
  \item \textbf{Modelar la incertesa d'un col·lectiu} (UD5): distribucions binomial i
        normal.
  \item \textbf{Inferir i criticar} (UD6): mostreig, intervals de confiança, significació i
        lectura crítica d'informes.
\end{enumerate}
\textbf{Els hàbits crítics adquirits:} mirar d'on surten les dades, expressar la incertesa,
no confondre correlació amb causa, no precipitar-se davant d'una diferència.
\end{idea}

\subsection{Exemples}

\begin{exemple}[una mateixa dada, mirada amb les quatre eines]
Davant d'un estudi que diu «el $60\%$ dels usuaris està satisfet»:
\begin{itemize}[leftmargin=1.4em, itemsep=2pt]
  \item \textbf{Descriure:} el $60\%$ és un paràmetre que resumeix la mostra.
  \item \textbf{Probabilitat/model:} descriu com es reparteix la satisfacció al col·lectiu.
  \item \textbf{Inferir:} el $60\%$ és una estimació de tota la població, amb un marge.
  \item \textbf{Criticar:} d'on surt la mostra? És representativa? Amb quin marge? Es pot
        generalitzar?
\end{itemize}
Una sola dada, i totes les eines del curs hi intervenen.
\end{exemple}

\begin{exemple}[per què és una competència professional]
Un tècnic de Serveis a la Comunitat haurà de \textbf{fonamentar decisions} en evidències:
llegir informes, encarregar estudis, avaluar intervencions. La mirada crítica amb dades li
permet no deixar-se enganyar per estudis mal fets o interessats, i defensar amb rigor les
necessitats dels col·lectius amb què treballa. No és només una competència acadèmica: és
\textbf{ciutadana i professional}.
\end{exemple}

\subsection{Exercicis resolts}

\begin{exresolt}[situar cada eina al seu pas]
Relaciona cada tasca amb el pas del recorregut: (a) calcular la mitjana d'unes edats;
(b) construir l'interval de confiança d'una enquesta; (c) valorar si una diferència entre
dos grups és casual.
\solucio
(a) \textbf{Descriure} (estadística descriptiva de 1r); (b) \textbf{Inferir} (intervals de
confiança, UD6); (c) \textbf{Inferir/criticar} (significació, UD6). Es veu com el curs
avança de descriure cap a inferir i criticar.
\resposta{(a) descriure; (b) inferir; (c) inferir/criticar.}
\end{exresolt}

\begin{exresolt}[l'hàbit crític en acció]
Un company diu: «aquest estudi ho demostra, ho posa aquí». Quins dos hàbits crítics li
recordaries abans de donar-ho per bo?
\solucio
Li recordaria: (1) \textbf{mirar d'on surten les dades} (quina mostra, representativa o no,
amb quin marge); i (2) \textbf{no confondre correlació amb causa} ni prendre una estimació
per una certesa. «Ho posa aquí» no és garantia: cal valorar-ne la solidesa.
\resposta{Mirar l'origen de les dades i no confondre estimació/correlació amb certesa/causa.}
\end{exresolt}

\subsection{Exercicis per fer}

\begin{exercicis}
\begin{exlist}
  \item Ordena els quatre passos del recorregut estadístic de l'etapa (descriure, mesurar
        la incertesa, modelar, inferir i criticar) i posa un exemple de cadascun.
  \item Per a cada tasca, digues a quin pas pertany: (a) fer un gràfic de barres; (b)
        calcular la probabilitat d'un cas; (c) llegir críticament un informe.
  \item Escriu els quatre hàbits crítics que t'emportes de la unitat i explica'n un amb un
        exemple.
  \item Per què diem que la mirada crítica amb dades és una competència «ciutadana», i no
        només acadèmica?
  \item \textbf{(Aplicació)} Pensa en una situació del teu futur professional (Serveis a la
        Comunitat) on hauràs de llegir o encarregar dades. Quines preguntes crítiques hi
        faràs?
  \item \textbf{(Reflexió)} Escriu una reflexió breu sobre què t'emportes d'aquest
        recorregut estadístic i com penses aplicar la mirada crítica amb dades.
\end{exlist}
\end{exercicis}

% ---------------------------------------------------------------------
\fullsolucions{Sessió 13}
\begin{solucions}
\begin{sollist}
  \item (1) Descriure —resumir dades amb paràmetres i gràfics; (2) mesurar la incertesa
        —probabilitat d'un cas; (3) modelar —distribucions binomial i normal; (4) inferir i
        criticar —mostreig, intervals, significació, lectura crítica. Exemples oberts.
  \item (a) descriure; (b) mesurar la incertesa (probabilitat); (c) inferir i criticar.
  \item Resposta oberta. Els quatre hàbits: mirar d'on surten les dades, expressar la
        incertesa, no confondre correlació amb causa, no precipitar-se davant d'una
        diferència.
  \item Resposta oberta. Idea: perquè qualsevol ciutadà rep contínuament dades (enquestes,
        titulars, estudis) i necessita interpretar-les i qüestionar-les per decidir i no
        deixar-se manipular; no és exclusiva de l'àmbit acadèmic.
  \item Resposta oberta. Idea: preguntes com d'on surten les dades, quina mostra i quin
        marge, si és representativa, si les diferències són significatives, quines
        limitacions té l'estudi.
  \item Resposta oberta.
\end{sollist}
\end{solucions}
